\subsection{特征标空间的分划}

命题 13. --- 设 A 是一个交换幺 Banach 代数。设 U₁ 与 U₂ 是 X(A) 的开子集，构成 X(A) 的一个分划。则 A 中存在唯一的幂等元 j，使得 Gelfand 变换 G(j) 在 U₁ 上等于 1，在 U₂ 上等于 0。

我们把空间 \(\mathsf{X}(\mathsf{A})\) 等同于 \(\mathbf{C}^{\mathrm{A}}\) 的一个紧子集，所借助的映射是 \(\chi \mapsto (\chi (a))_{a\in \mathrm{A}}\)（参见 I, p. 6, 第 6 小节以及 I, p. 29 定理 1 的推论）。子集 \(\mathrm{U}_1\) 与 \(\mathrm{U}_2\)（属于一致空间 \(\mathbf{C}^{\mathrm{A}}\)）是紧的且不相交。

由 TG, II, p.~31 的命题 4，存在 A 的有限子集 M 以及不相交的开集 \(\mathrm{V}_{1}\) 与 \(\mathrm{V}_{2}\)（属于 \(\mathbf{C}^{\mathrm{M}}\)），使得

\[
p (\mathrm{U} _ {1}) \subset \mathrm{V} _ {1}, \qquad p (\mathrm{U} _ {2}) \subset \mathrm{V} _ {2},
\]

其中 \(p\) 是 \(\mathbf{C}^{\mathrm{A}}\) 到 \(\mathbf{C}^{\mathrm{M}}\) 上的典范投影。

设 \(a_1, \ldots, a_n\) 是 M 的互不相同的元素，并把 \(\mathbf{C}^{\mathrm{M}}\) 等同于 \(\mathbf{C}^n\)。我们有 \(\mathrm{Sp}_{\mathrm{A}}^{n}(a_{1}, \ldots, a_{n}) \subset p(\mathsf{X}(\mathrm{A})) \subset \mathrm{V}_{1} \cup \mathrm{V}_{2}\)，因为 \(\mathrm{U}_{1} \cup \mathrm{U}_{2} = \mathsf{X}(\mathrm{A})\)。设 \(f\) 是 \(\mathrm{V}_{1} \cup \mathrm{V}_{2}\) 上的函数，在 \(\mathrm{V}_{1}\) 上等于 1，在 \(\mathrm{V}_{2}\) 上等于 0。我们有 \(f \in \mathcal{O}(\mathrm{V}_{1} \cup \mathrm{V}_{2})\)。令 \(j = f(a_{1}, \ldots, a_{n})\)。由于 \(f^{2} = f\)，有 \(j^{2} = j\)。由 I, p.~66 的推论 1，我们有 \(\chi(j) = 1\)（若 \(\chi \in \mathrm{U}_{1}\)），而 \(\chi(j) = 0\)（若 \(\chi \in \mathrm{U}_{2}\)），这就证明了所需幂等元的存在性。

另一方面，若 \(j_{1}\) 是 A 的幂等元，关系式 \(j^{2} = j\) 与 \(j_{1}^{2} = j_{1}\) 蕴涵 \((j - j_{1})(j + j_{1} - 1) = 0\)。若 \(\mathcal{G}(j_{1}) = \mathcal{G}(j)\)，则 \(j + j_{1} - 1\) 的 Gelfand 变换取值于 \(\{-1, 1\}\)，因而 \(j + j_{1} - 1\) 可逆（I, p.~37 的命题 6），于是 \(j = j_{1}\)。

推论. --- 设 A 是一个交换幺 Banach 代数。下列断言等价：

\begin{enumerate}
\def\labelenumi{\alph{enumi})}
\item
  特征标空间 X(A) 不是连通的；
\item
  A 存在异于 0 与 1 的幂等元。
\item
  代数 A 同构于两个非零 Banach 代数的乘积。
\end{enumerate}

该命题表明 \(a\)) 蕴涵 \(b\))。若 \(j\) 是 A 的幂等元，令 \(\mathrm{I}_1 = j\mathrm{A}\)，\(\mathrm{I}_2 = (1 - j)\mathrm{A}\)。那么 \(\mathrm{I}_1\) 与 \(\mathrm{I}_2\) 是 A 的闭理想，且 \(\mathrm{I}_1 + \mathrm{I}_2 = \mathrm{A}\)。若 \(j \notin \{0,1\}\)，则理想 \(\mathrm{I}_1\) 与 \(\mathrm{I}_2\) 都异于 A。另一方面，理想 \(\mathrm{I}_1\)（相应地，\(\mathrm{I}_2\)）是由 A 的元素 \(x\)（满足 \(jx = x\)，相应地，\((1 - j)x = x\)）组成的集合，因而 \(\mathrm{I}_1 \cap \mathrm{I}_2 = \{0\}\)。于是代数 A 等同于乘积 \(\mathrm{A} / \mathrm{I}_1 \times \mathrm{A} / \mathrm{I}_2\)。因此断言 \(b\)) 蕴涵 \(c\))。最后，若 A 同构于 \(\mathrm{A}_1 \times \mathrm{A}_2\)，则空间 \(\mathsf{X}(\mathsf{A})\) 等同于 \(\mathsf{X}(\mathsf{A}_1)\) 与 \(\mathsf{X}(\mathsf{A}_2)\) 的和空间（I, p.~6, 第 6 小节），因而 \(c\)) 蕴涵 \(a\))。

命题 14. --- 设 A 是一个无根的交换 Banach 代数。A 具有单位元，当且仅当 X(A) 是紧的。

条件是必要的（I, p.~29, 推论）。设 \(\mathsf{X}(\mathsf{A})\) 紧。设 \(\widetilde{\mathsf{A}}\) 是由 A 添加单位元而得到的 Banach 代数，并把 \(\mathsf{X}'(\mathsf{A})\) 等同于 \(\mathsf{X}(\widetilde{\mathsf{A}})\)。\(\mathsf{X}(\mathsf{A})\) 在 \(\mathsf{X}(\widetilde{\mathsf{A}})\) 中的余集化为一个特征标 \(\chi_0\)：它是 \(\widetilde{\mathsf{A}}\) 的特征标，其核为 A。集合 \(\mathsf{X}(\mathsf{A})\) 与 \(\{\chi_0\}\) 在 \(\mathsf{X}(\widetilde{\mathsf{A}})\) 中是开的。由命题 13，存在元素 \(j \in \mathsf{A}\)，使得 \(\chi(j) = 1\)（对 \(\chi \in \mathsf{X}(\mathsf{A})\)），且 \(\chi_0(j) = 0\)。于是我们有 \(j \in \mathsf{A}\)。

然后设 A 中的 \(x\)。有 \(\chi(jx) = \chi(x)\)（对所有 \(\chi \in \mathsf{X}(\mathsf{A})\)），由于 A 无根，故 \(jx = x\)。于是 \(j\) 是 A 的单位元。

命题 15. --- 设 A 是一个交换 Banach 代数，\(I_{1}\) 是 A 的理想，\(F_{1}\) 是由 \(\chi \in \mathsf{X}(A)\)（在 \(I_{1}\) 上取值为零）组成的集合。设 \(F_{2}\) 是 X(A) 中与 \(F_{1}\) 不相交的子集，对 Jacobson 拓扑是闭的，且对弱拓扑是紧的。则存在 \(u \in \text{I}_{1}\)，使得 \(\mathcal{G}(u) = 1\) 在 \(F_{2}\) 上成立。

设 \(I_{2}\) 是属于 \(F_{2}\) 的特征标的核的交。Banach 代数 A/ \(I_{2}\) 无根（I, p.~38 的命题 8）。由于 \(F_{2}\) 对 Jacobson 拓扑是闭的，X(A) 中在 \(I_{2}\) 上取值为零的元素只有 \(F_{2}\) 中的那些（参见 I, p.~13）。因此，赋以 X(A) 的弱拓扑所诱导的拓扑的 \(F_{2}\)，等同于赋以弱拓扑的 X(A/ \(I_{2}\) )（I, p.~9, 第 7 小节）。由于 \(F_{2}\) 弱紧，代数 A/ \(I_{2}\) 具有单位元（命题 14）。

于是我们有 \(I_{1} + I_{2} = A\)。事实上，若不然，\((I_{1} + I_{2})/I_{2}\) 将是真理想，从而包含在 \(A/I_{2}\) 的某个非零特征标的核中（I, p.~30, 定理 2）。这个特征标通过与典范投影 \(A \rightarrow A/I_{2}\) 复合，将定义 A 的一个非零特征标 \(\chi\)，它在 \(I_{1}\) 与 \(I_{2}\) 上取值为零，因而属于 \(F_{1} \cap F_{2}\)，与假设矛盾。

由于 \(I_{1} + I_{2} = A\)，存在 \(u \in I_{1}\)，其类在 \(A/I_{2}\) 中是 \(A/I_{2}\) 的单位元。那么 \(\chi(u) = 1\)（对所有 \(\chi \in F_{2}\)），证明完毕。

推论. --- 设 A 是一个交换 Banach 代数。设 \(F_{1}\) 与 \(F_{2}\) 是 \(\mathsf{X}(\mathsf{A})\) 的两个对 Jacobson 拓扑为闭的不相交子集。假设 \(F_{2}\) 弱紧。则存在 \(u \in A\)，使得 \(\mathcal{G}(u) = 1\) 在 \(F_{2}\) 上成立，且 \(\mathcal{G}(u) = 0\) 在 \(F_{1}\) 上成立。

